\section{怎样判断参数好不好？从损失函数到梯度下降}
\label{intro:16-loss-gradient}

\begin{idea}
\noindent\emph{怎样根据误差调整参数？}\par

上一章试着猜了带参数的函数，但还没有统一的尺度判断它猜得如何。
先定义损失，再寻找怎样沿着损失降低的方向改变参数。
\end{idea}

\subsection{预测值、真实值、误差与损失}
\phantomsection\label{concept:loss}

设一个模型 $\widehat y=wx$，先不考虑偏置。样本为：
\[
(x_1,y_1)=(1,1),\quad(x_2,y_2)=(2,2).
\]
如果猜 $w=0$，两个预测都是零；如果猜 $w=1$，两个预测恰好是 $1,2$。为了量化差距，定义\textbf{均方误差（MSE）}：
\begin{equation}
L(w)=\frac1n\sum_{i=1}^{n}(\widehat y_i-y_i)^2.
\end{equation}
其中 $n$ 是样本数，$L$ 是\textbf{损失（loss）}。平方避免正负误差相互抵消，还使误差较大的样本被更多惩罚。

把两个样本代入：
\[\begin{aligned}
L(w)&=\tfrac12[(w-1)^2+(2w-2)^2]\\
&=\tfrac52(w-1)^2.
\end{aligned}\]
于是 $L(0)=2.5$、$L(0.5)=0.625$、$L(1)=0$。
训练的目标变成寻找让损失尽可能小的参数。

\subsection{先用导数直接解：为什么还需要梯度下降？}

现在先不要急着迭代。既然上面的损失是一个开口向上的二次函数，
我们完全可以用熟悉的求极值办法：\textbf{对参数求导，令导数等于零}。
导数描述参数改变一点时损失的局部变化趋势，计算得到
\[
\frac{\mathrm dL}{\mathrm dw}=5(w-1)=0
\quad\Longrightarrow\quad \boxed{w=1}.
\]
这是这道简单题的\textbf{全局最小值}：代回去，两个样本的误差都变为零。
既然直接算出答案了，为什么后面还要学“一小步一小步地调”？

因为\textbf{简单回归能直接解，不代表复杂网络也能直接解}。
多层神经网络的损失同时依赖大量参数和非线性运算；即使写出
“所有偏导数为零”的方程，通常也不能像这条抛物线一样直接求出可用的最优参数。
而且导数为零只表示驻点，并不总是最小值。

因此换一种更可实施的目标：\textbf{不先求整个问题的最终答案，
只算当前位置朝哪个方向调参数，更可能让损失下降。}
这才需要梯度下降。它是数值优化方法，不是“求导失效了”。

\subsection{导数是什么意思？先用一个小改动体验}

\noindent\emph{数学随手查：}导数表示函数在某个位置的局部变化率，相关符号见\hyperref[app:math-tools]{数学速查}。


继续使用刚才求出的 $L(w)=\tfrac52(w-1)^2$，目前 $w=0$。稍微把 $w$ 增大，损失会减小；而从 $w=2$ 开始增大，损失反而变大。导数衡量函数在某个点\textbf{局部怎样变化}：
\begin{equation}
\frac{\mathrm dL}{\mathrm dw}
=\lim_{\varepsilon\to0}
\frac{L(w+\varepsilon)-L(w)}{\varepsilon}.
\end{equation}
这个比值比较参数发生微小变化时，损失相应改变多少；取极限便得到该点的瞬时变化率。

本例中，
\[
L(w)=\tfrac52(w-1)^2,\qquad
\frac{\mathrm dL}{\mathrm dw}=5(w-1).
\]
在 $w=0$ 时导数是 $-5$：增大 $w$ 有望使损失下降；在 $w=2$ 时导数是 $5$：应该减小 $w$。

\subsection{梯度下降：往损失增大的反方向走}
\phantomsection\label{concept:gradient-descent}

\begin{equation}
\boxed{w_{\mathrm{new}}=
w_{\mathrm{old}}-\eta\frac{\mathrm dL}{\mathrm dw}}
\label{eq:gradient-simple}
\end{equation}
其中 $\eta$ 是\textbf{学习率}（learning rate），控制每次修改的幅度。

假设 $w_0=0$、$\eta=0.1$。第一次更新：
\[
w_1=0-0.1\times(-5)=0.5.
\]
第二次，导数是 $5(0.5-1)=-2.5$：
\[
w_2=0.5-0.1\times(-2.5)=0.75.
\]
损失也从 $2.5$ 变为 $0.625$，再变为 $0.15625$。这就是\textbf{迭代优化}：反复评估局部方向，逐渐逼近低损失区域。

% Local probes of the same loss at b=0; no separate training trajectory.
\par\noindent\begin{minipage}{\linewidth}
\tagstructbegin{tag=Figure,alttext={固定偏置 b 为零，横轴是斜率参数 w，纵轴是损失。三个试算点分别为零点九对应一点七四、一对应三分之五、一点一对应约一点六零六七。曲线在当前点向右下方延伸，虚线切线的斜率为负三分之二，因此在附近增大 w 能减小损失。}}
\begin{center}
\begin{tikzpicture}
\begin{axis}[
 width=112mm,height=52mm,
 xmin=.78,xmax=1.24,ymin=1.48,ymax=1.91,
 axis lines=left,xlabel={$w$},ylabel={$L(w,0)$},
 xtick={.9,1,1.1},ytick={1.606667,1.666667,1.74},
 yticklabels={$1.6067$,$5/3$,$1.74$},
 tick label style={font=\footnotesize},label style={font=\small},
 legend style={at={(.5,-.30)},anchor=north,draw=none,font=\footnotesize,legend columns=2},
]
\addplot[flowcolor,very thick,domain=.8:1.22,samples=60] {((1-x)^2+4+(x-2)^2)/3};
\addlegendentry{损失曲线}
\addplot[residualcolor,thick,dashed,domain=.8:1.22] {5/3-2/3*(x-1)};
\addlegendentry{当前点的切线}
\addplot[datacolor,only marks,mark=*,mark size=2.4pt,forget plot]
 coordinates {(.9,1.74) (1,1.666667) (1.1,1.606667)};
\draw[-{Latex},flowcolor,thick] (axis cs:.94,1.855)--(axis cs:1.15,1.855)
 node[midway,above,font=\footnotesize] {增大 $w$};
\node[font=\footnotesize,anchor=west] at (axis cs:1.025,1.745) {切线斜率 $-2/3$};
\end{axis}
\end{tikzpicture}
\captionof{figure}{固定 $b=0$ 的局部观察：试算结果与负斜率指向同一下降方向}
\label{fig:loss-descent}
\end{center}
\tagstructend
\end{minipage}\par


\subsection{同时调整多个参数：偏导数与梯度}
\phantomsection\label{concept:gradient}

假设模型是 $\widehat y=wx+b$，参数是 $w,b$ 两个数。损失写成 $L(w,b)$：
\[
\nabla L=
\left(\frac{\partial L}{\partial w},\frac{\partial L}{\partial b}\right),
\qquad
(w,b)_{\mathrm{new}}=(w,b)_{\mathrm{old}}-\eta\nabla L.
\]
\textbf{偏导数}就是求某一个参数方向的变化率时，暂时把其他参数视为常数。把两个偏导数排成一组，得到\textbf{梯度}。以后参数增多时再用向量和矩阵表示。

\begin{warningbox}
导数等于零不一定就是全局最小值；还可能是最大值或鞍点。神经网络损失面一般复杂且不凸，梯度下降也没有保证每次都找到全局最优。学习率太大可能震荡或发散，太小则训练缓慢。
\end{warningbox}

\subsection{一次训练要做哪四件事？}
\begin{enumerate}
\item \textbf{前向传播}：用当前权重计算 $\widehat y$。
\item \textbf{计算损失}：比较 $\widehat y$ 和实际目标 $y$。
\item \textbf{求梯度}：判断每个权重应该朝哪个方向修改。
\item \textbf{更新参数}：执行 $\theta\leftarrow\theta-\eta\nabla_\theta L$。
\end{enumerate}

\par\medskip\noindent{\sffamily\footnotesize\color{blue}\textbf{读完自测}}\quad\phantomsection\label{check:16-loss-gradient}$L(w)=(w-1)^2$ 在 $w=0$ 处的导数是多少？\quad{\footnotesize\hyperref[answer:16-loss-gradient]{核对解答。}}\par\medskip

\begin{conclusion}
损失函数把“模型好不好”变成数字；梯度提供局部下降方向；学习率控制步长。下一章先问：即使参数可以优化，若函数本身只能画直线，又怎样表示弯曲的关系？
\end{conclusion}
